\scanpage{31}
% L1–2
6. \uline{Caractérisation des modules fondamentaux d'anneaux Coh. Mac.} \uline{Anneaux de Gorenstein}.

% L3–4
\deleted{\uline{Théorème 5} Soient $A$ un anneau local noeth, $M$ un $A$-module}

% L5–6, large cross over the first version
\deleted{\uline{Prop. 13} Soient $A$ un anneau local noeth\inserted{ de corps rés. $k$,} $M$ un $A$-module, $n$ un entier $>0$, on suppose}
% L7
\[
\deleted{\operatorname{Ext}^{n-1}(k,M)=0}
\]
% L8–11
\deleted{Alors les conditions suivantes sont équivalentes}

\deleted{a) $\operatorname{Ext}^{n}(k,M)\simeq k$ et $\operatorname{Ext}^{n+1}(k,M)=0$}

\deleted{b) $N\longrightarrow\operatorname{Ext}^{n}(N,M)$ est dualisant}

\deleted{c) $H^{n}(M)$ est un module dualisant}

% L12–13
\deleted{On a alors $n=\dim A$, les idéaux pr. associés à $A$ sont \unsure{tous de dim $n$|tous de la dim.|étant de dim. $n$}, et $\lg_{\mathfrak p}(M)=\lg_{\mathfrak p}(A)$ pour tt tel $\mathfrak p$}

% L14–19, moved clauses connected by lines
\deleted{\uline{Cor. 1} Soit $\mathfrak p$ un idéal pr. minimal de $A$, les conditions précédentes équivalent alors} \deleted{à la conjonction des} \deleted{aux conditions : $n=\dim A$,} \inserted{\deleted{Supposons $n=\dim A$.}}

\deleted{$0<\lg_{\mathfrak p}M\leq\lg_{\mathfrak p}A$ et à une \unsure{conditions}} \deleted{$\unsure{\operatorname{Ext}^{n-1}(k,M)\ne0|\operatorname{Ext}^{n-1}(k,M)=0}$, $n=\dim A$,} \inserted{\deleted{et $0<\lg_{\mathfrak p}M\leq\lg_{\mathfrak p}A$}}

\deleted{et à l'une des conditions équivalentes suivantes}

% L20–22
\deleted{a') $\operatorname{Ext}^{n+1}(k,M)=0$}

\deleted{b') $N\longrightarrow\operatorname{Ext}^{n}(N,M)$ est exact}

\deleted{c') $H^{n}(M)$ est injectif} $\downarrow$

% L23–26, fresh version at foot
\deleted{\uline{Lemme 7}} \inserted{\uline{Proposition 13}} Soit $A$ anneau local\inserted{ de dim $n$}, $M$ un $A$-module, \deleted{\illegible{1}}

\deleted{\illegible{2}} Supposons $\boxed{\operatorname{Ext}^{n-1}_{A}(k,M)=0}$, \inserted{on a $\boxed{\operatorname{Ext}^{n}(N,M)\simeq\operatorname{Hom}(N,H^{n}(M))}$ pour $N$ de long. finie} \deleted{Alors} et les conditions suivantes sont équivalentes

% L27–29
a') $\operatorname{Ext}^{n+1}(k,M)=0$

b') $N\longrightarrow\operatorname{Ext}^{n}(N,M)$ est exact ($N$ de long finie)

c') $H^{n}(M)$ est injectif

% archive sheet number: 30
